\section{A uniform bound on every actual tensor orbit}
\label{sec:general-tensor-orbit-bound}

\paragraph{Status and provenance.}
D-TD16, reserved question Q2ef, is a separate derivation from the
root-proposed orbit budget, with the complete proof and exact constants
checked by \texttt{/root/imprimitive}. The strategy and its four
factor-fixing entries are shared root contributions, not an independently
originated argument. The proof retains the entire external partner and
uses the accepted complete decomposition D-TD11--13, the actual-lift
trace and quarter bounds D-TD14, and the endpoint-chord idea of HD-TC01.
Complete author, root and independent whole-proof readings pass without mathematical repair; their exact source and computation scopes are recorded separately.
It changes none of those earlier statements.

For a semisimple operator \(T\) on a space of dimension \(e\), put
\(\nu(T)=e-\max_z\dim\ker(T-zI)\). Actual moving rank
\(\operatorname{rank}(1-T)\) is at least \(\nu(T)\).

\begin{theorem}[D-TD16]
Let \(J\subset\operatorname{GL}(V)\) be a finite primitive linear
group, with \(\dim V=e>1\). Suppose an actual connected
\(J\)-Galois cover of \(\mathbf P^1\) has genuine generating
product-one inertia \(g_1,\ldots,g_r\), where
\[
 r\ge12,\qquad g_i\ne1,\qquad
 \sum_{i=1}^r\operatorname{rank}(1-g_i\mid V)=2e+6.
 \tag{TOR.1}
\]
Perform one complete recursive decomposition of the projective image
of \(J\) as in D-TD12. Let \(m_*\) be the largest alternating
natural degree among the quasisimple factors at every recursion layer,
with \(m_*=1\) if there is no alternating factor, and put
\[
                         M=\max\{80,m_*-1\}.
 \tag{TOR.2}
\]
Consider any actual transitive orbit of \(t\) full tensor factor
algebras normalized and permuted by \(J\), of common local dimension
\(a\ge2\), with their entire commutant retained:
\[
 V=(W_1\otimes\cdots\otimes W_t)\otimes R,
 \qquad \dim W_i=a,\qquad \dim R\ge1.
 \tag{TOR.3}
\]
Then
\[
                            t\le\lfloor M/2\rfloor.
 \tag{TOR.4}
\]
In particular, this holds simultaneously for every orbit in the chosen
complete recursive decomposition. No assumption on the transitive
permutation image, its number of genuine projected marks, or the
dimension or action of \(R\) is required.

If \(m_*=n\ge81\), the combined image of \(J\) on all factor
orbits of that decomposition has no section isomorphic to \(A_n\)
or \(S_n\). Here a section means a quotient of a subgroup.
\end{theorem}

\begin{proof}
The projective image of a finite primitive linear group is primitive.
Apply D-TD13 to the whole space \(V\), with all layers included in
\(m_*\). Every nonscalar operator in \(J\) satisfies
\(\nu(g)/e\ge1/M\). A scalar genuine inertia entry is \(zI\)
with \(z\ne1\), so its actual moving rank is \(e\). Thus
\[
             \operatorname{rank}(1-g_i\mid V)\ge e/M
                         \quad\hbox{for every }i.
 \tag{TOR.5}
\]
This is a bound on the entire coefficient representation. It also
charges entries acting only on \(R\), or only through a scalar;
no factor or multiplicity space has been omitted.

Conjugation on the factor algebras in (TOR.3) gives an actual
homomorphism \(x:J\to S_t\) with transitive image. Their product
algebra has commutant \(\operatorname{End}(R)\), which is preserved.
Undoing the permutation and then the inner automorphisms of the full
factor algebras writes each operator as a weighted tensor permutation
times an operator on \(R\). This is precisely the actual algebra
normalizer used in D-TD14; it does not select stable individual
representations \(W_i\).

We record the chord bound with the full lift order. For \(T\in J\)
let \(\pi=x(T)\), let \(L\) be the order of \(T\), and let
\(f_j\) be the number of fixed letters of \(\pi^j\). By D-TD14,
including its arbitrary external operator and all kernel powers,
\[
 \frac{|\operatorname{tr}(T^j)|}{e}
 \le a^{-\operatorname{ind}(\pi^j)}
 \le a^{-(t-f_j)/2}
 \le a^{-t/2}+(1-a^{-t/2})\frac{f_j}{t}.
 \tag{TOR.6}
\]
The second inequality uses that every nontrivial permutation cycle
has length at least two. The third is the endpoint chord of the
convex exponential function on \([0,t]\). Since the order of
\(\pi\) divides \(L\), the average of \(f_j\), for
\(0\le j<L\), is the number of cycles of \(\pi\). Fourier
inversion over this full \(L\), applied to any eigenvalue of \(T\),
therefore yields
\[
 \frac{\nu(T)}e\ge
           (1-a^{-t/2})\frac{\operatorname{ind}(\pi)}{t}.
 \tag{TOR.7}
\]
In particular, no power in the kernel of \(x\) has been assumed
scalar. Independently, D-TD14 gives
\[
                 \pi\ne1\quad\Longrightarrow\quad\nu(T)/e\ge1/4.
 \tag{TOR.8}
\]
Both inequalities hold for every target eigenvalue, so scalar choices
in the local tensor maps do not change them.

Suppose, to obtain a contradiction, that \(t>\lfloor M/2\rfloor\).
As \(M\ge80\) is an integer, \(t\ge41\) and
\(e\ge a^t\ge2^t>24\). Let \(r_0\) count the original entries
with \(x(g_i)\ne1\). From (TOR.1) and (TOR.8),
\[
                         r_0/4\le2+6/e<9/4,
                         \qquad r_0\le8.
 \tag{TOR.9}
\]
Consequently at least four of the original genuine entries fix
every factor algebra in the chosen orbit. By (TOR.5) their total
moving rank is at least \(4e/M\). These entries may act nontrivially
only through the external partner; their cost is still retained.

The actual quotient by \(\ker x\), followed by a letter stabilizer,
is a connected degree-\(t\) cover of \(\mathbf P^1\). Its branch
tuple is \(x(g_1),\ldots,x(g_r)\), with identity images omitted.
Riemann--Hurwitz, with nonnegative quotient genus, gives
\[
                  \sum_i\operatorname{ind}(x(g_i))\ge2t-2.
 \tag{TOR.10}
\]
Apply (TOR.7) to the moving entries and (TOR.5) to four of the
factor-fixing entries. These are disjoint sets of entries. Equations
(TOR.1) and (TOR.10) imply
\[
 2+6/e\ge(1-a^{-t/2})(2-2/t)+4/M,
 \qquad
 \frac4M-\frac2t
       \le2\,2^{-t/2}+6\,2^{-t}.
 \tag{TOR.11}
\]
For the second inequality we used \(a\ge2\), \(e\ge2^t\),
and \(1-1/t\le1\).

On the other hand \(M\le2t-1\), so
\[
 \frac4M-\frac2t\ge
       \frac2{t(2t-1)}>\frac1{t^2}.
 \tag{TOR.12}
\]
The following two exact inequalities hold for all integers \(t\ge32\):
\[
                         2^t>16t^4,
                         \qquad 2^t>12t^2.
 \tag{TOR.13}
\]
Their base values at \(32\) are immediate integer comparisons.
For the induction, \((1+1/t)^4\le(33/32)^4<2\), and
\((1+1/t)^2\le(33/32)^2<2\). Thus each right-hand side grows
by a factor strictly less than two at every next integer.
The first inequality in (TOR.13), after taking its positive square
root, and the second inequality give respectively
\[
                  2\,2^{-t/2}<\frac1{2t^2},
                  \qquad6\,2^{-t}<\frac1{2t^2}.
 \tag{TOR.14}
\]
This contradicts (TOR.11)--(TOR.12), since \(t\ge41\).
Hence (TOR.4) holds.

It remains to prove the simultaneous permutation assertion. Suppose
\(m_*=n\ge81\). Then \(M=n-1\). Let \(t_1,\ldots,t_s\)
be all orbit lengths in the complete decomposition, including any
one-letter orbits, and write its combined permutation image as
\[
                  P\subset\prod_{i=1}^s S_{t_i},
             \qquad t_i\le\lfloor(n-1)/2\rfloor<n.
 \tag{TOR.15}
\]
Each \(S_{t_i}\) has order at most \((n-1)!<n!/2=|A_n|\).
The product order itself need not be small, so a further argument
is necessary. If a subgroup \(H\le P\) had a surjection
\(\varphi:H\to A_n\), let \(H_i\) be the kernel on \(H\)
of projection to its first \(i\) coordinates. Thus \(H_0=H\),
\(H_s=1\), and all \(H_i\) are normal in \(H\). Simplicity of
\(A_n\) makes each \(\varphi(H_i)\) either \(A_n\) or one.
At the first drop, \(A_n\) is a quotient of
\(H_{i-1}/H_i\). Projection to the \(i\)-th coordinate embeds
this latter group in \(S_{t_i}\), contradicting its smaller order.
Thus there is no \(A_n\) section. An \(S_n\) section would,
by taking the inverse image of its subgroup \(A_n\), give an
\(A_n\) section as well, and is equally impossible.
\end{proof}

\paragraph{Scope of the reduction.}
The use of twelve marks is only at the original faithful coefficient
level in (TOR.1); it is not imposed on any projected factor tuple.
The actual transitive permutation quotient in (TOR.10) may be any
finite group. The integer \(m_*\) must include every recursive layer
of the whole coefficient representation. A first-core maximum or an
incomplete external-partner decomposition cannot be substituted.

The conclusion bounds permutation-orbit lengths in terms of \(m_*\),
not \(m_*\) itself. It does not exclude alternating quotients arising
elsewhere in the actual group, classify residual Heisenberg or
quasisimple carriers, prove a bound on original inducing indices,
or solve any coefficient-field or polarization descent problem.
No new primary classification or character-table input enters here;
the only non-elementary support interface is exactly the already
reviewed D-TD13, with its accepted HD-FT20--21 and D-TD08 dependencies.
